%!TEX root = ../chapter02.tex 
%----------------------------------------------------------------------------------------
%	
%----------------------------------------------------------------------------------------

%----------------------------------------------------------------------------------------
%	 
%---------------------------------------------------------------------------------------- 
\newpage 
\loadgeometry{widelayout}  
\pagestyle{scrheadings} % Print headers again  
\KOMAoptions{
	headwidth=textwithmarginpar,
	headsepline=true,
	headsepline= 0.5pt,
	footwidth=textwithmarginpar,
	paper=A4,
	paper=portrait,
}     

\onehalfspacing    
\section{Factorisation à l'aide d'identités remarquables}

\begin{exercise} Compléter les égalités suivantes par des identités remarquables :
	\begin{align*}
		\big(x-\ldots\ldots\big)^2&=\ldots\ldots x^2-\ldots\ldots x+36 \rule{0pt}{1.6em} \\
		\big(\ldots\ldots+\ldots\ldots\big)^2&=9  x^2+\ldots\ldots x+16 \rule{0pt}{1.6em} \\ 
		\big(\ldots\ldots x+\ldots\ldots\big)^2&=\ldots\ldots+ 20 x+4\rule{0pt}{1.6em} \\
		\big(\ldots\ldots x+\ldots\ldots\big)\big(\ldots\ldots - \ldots\ldots\big)&=4x^2 -36\rule{0pt}{1.6em} \\ 
		\big(\ldots\ldots x+\ldots\ldots\big)\big(\ldots\ldots - \ldots\ldots\big)&=36x^2 -4\rule{0pt}{1.6em} 
	\end{align*} 
\end{exercise}

\begin{example}[factoriser des expressions sous la forme $a^2\pm 2ab+b^2$]\hfill  
	
\noindent\parbox{0.45\linewidth}{	
$\begin{WithArrows}
A& = x^2 +8x +16 \rule{0pt}{1.6em} \\
&= \big ( \qquad \big)^2 \ldots 2\times\big ( \qquad \big)\times \big ( \qquad \big) \ldots \big ( \qquad \big)^2  \rule{0pt}{1.6em} \\
& =\big ( \qquad\qquad\qquad \big)^2\rule{0pt}{1.6em}   
\end{WithArrows}$ 	
}\hfill\parbox{0.45\linewidth}{	
$\begin{WithArrows}
	B& = x^2 -6x +9 \rule{0pt}{1.6em} \\
	&= \big ( \qquad \big)^2 \ldots 2\times\big ( \qquad \big)\times \big ( \qquad \big) \ldots \big ( \qquad \big)^2  \rule{0pt}{1.6em} \\
	& =\big ( \qquad\qquad\qquad \big)^2\rule{0pt}{1.6em}   
\end{WithArrows}$ 	
}
	
\noindent\parbox{0.45\linewidth}{	
	$\begin{WithArrows}
		C& = 25x^2 +80x +64 \rule{0pt}{1.6em} \\
		&= \big ( \qquad \big)^2 \ldots 2\times\big ( \qquad \big)\times \big ( \qquad \big) \ldots \big ( \qquad \big)^2  \rule{0pt}{1.6em} \\
		& =\big ( \qquad\qquad\qquad \big)^2\rule{0pt}{1.6em}   
	\end{WithArrows}$ 	
}\hfill\parbox{0.45\linewidth}{	
	$\begin{WithArrows}
		D& = \frac{1}{4}x^2+x+1 \rule{0pt}{1.6em} \\
		&= \big ( \qquad \big)^2 \ldots 2\times\big ( \qquad \big)\times \big ( \qquad \big) \ldots \big ( \qquad \big)^2  \rule{0pt}{1.6em} \\
		& =\big ( \qquad\qquad\qquad \big)^2\rule{0pt}{1.6em}   
	\end{WithArrows}$ 	
}	 
\end{example} 
\begin{exercise}\label{chap02:exo:factoriser:identite:niveau1} Pour $x\in\mathbb{R}$. Factoriser au maximum les expressions suivantes à l'aide d'identités remarquables :
	\begin{multicols}{3}
		\begin{enumerate}[label={$\Alph*(x)=$},left= 0pt]
			\item $x^2 + 2x + 1$ 
			\item $1 - 2x + x^2$ 
			\item $x^2 + 6x + 9$
			\item $4x^2 + 4x + 1$ 
			\item $4x^2 + 12x + 9$ 
			\item $4x^2 - 12x + 9$
			\item $x^2 + x + \tfrac{1}{4}$  
			\item $9x^2-6\sqrt{5}x + 5$  
			\item $3x^2-2\sqrt{3}x + 1$  
		\end{enumerate}
	\end{multicols}
\end{exercise}
\begin{exercise}\label{chap02:exo:factoriser:identite:niveau2} \emoji{cactus} Pour $x\in\mathbb{R}$. Factoriser au maximum  :
	\begin{multicols}{3}
		\begin{enumerate}[label={$\Alph*(x)=$},left= 0pt]
			\item $x^2-36$ 
			\item $4x^2-36$
			\item $36x^2-4$ 
			\item $(2x-1)^2-25$
			\item $(3x+7)^2-16$
			\item $4(2x-1)^2-100$  
			\item $9(2x-1)^2-4x^2$
			\item $4(2x-1)^2-25x^2$  
			\item $4(2x-1)^2-25(x+3)^2$  
		\end{enumerate}
	\end{multicols}
\end{exercise}
\begin{proof}[solution de l'exercice \ref{chap02:exo:factoriser:identite:niveau1} ]  
	\begin{pycode}
from re import sub
import sympy as sp

x, y, z, t = sp.symbols('x y z t')
k, m, n = sp.symbols('k m n', integer=True)
f, g, h = sp.symbols('f g h', cls=sp.Function)

# Create a list of functions to include in the table
funcsfac = [ 
'x*x+2*x+1',
'1-2*x+x**2',
'x*x+6*x+9',
'4*x**2+4*x+1',
'4*x**2+12*x+9',
'4*x**2-12*x+9',
'x**2+x+0.25'  
]
n=0
L=['A','B','C','D','E','F', 'G', 'H', 'I']
for func in funcsfac :  
	myfactor =  eval('sp.factor(func)')
	print( '$'+ str(L[n])  + "=" + sp.latex(myfactor) + '$; ' ) 
	n     = n + 1
\end{pycode}
$H=(3x-5)^2$; $I=(\sqrt{3} x-1)^2$.
\end{proof}
\begin{proof}[solution de l'exercice \ref{chap02:exo:factoriser:identite:niveau2} ]  
	\begin{pycode}
from re import sub
import sympy as sp

x, y, z, t = sp.symbols('x y z t')
k, m, n = sp.symbols('k m n', integer=True)
f, g, h = sp.symbols('f g h', cls=sp.Function)

# Create a list of functions to include in the table
funcsfac = [
'x**2-36',
'4*x**2-36',
'36*x**2-4',
'(2*x-1)**2-25',
'(3*x+7)**2-16',
'4*(2*x-1)**2-100',
'9*(2*x-1)**2-4*x**2',
'4*(2*x-1)**2-25*x**2',
'4*(2*x-1)**2-25*(x+3)**2'
]
n=0
L=['A','B','C','D','E','F', 'G', 'H', 'I']
for func in funcsfac :
	myfactor =  eval('sp.factor(func)')
	print( '$'+ str(L[n])  + "=" + sp.latex(myfactor) + '$; ' )
	n     = n + 1
	\end{pycode}
\end{proof}

%----------------------------------------------------------------------------------------
%	
%----------------------------------------------------------------------------------------




